\section{What Has Actually Been Defined}\label{sec:defined}

Before asking why there are two judgments it is worth establishing, from the acts themselves, that there are two---and exactly what about each of them the Church has bound herself to. The relevant texts are few, short, and unusually precise. They are also, read in sequence, more interesting than their reputation: three of them assert the immediacy of a man's destiny at death and, in the same sentence, retain the general judgment. The tradition did not stumble into the tension this article is about. It wrote the tension into its definitions on purpose.

\subsection{One slot on the list, two events in the acts}

Part of the difficulty is manufactured by the vehicle that carries the doctrine to most people. ``The four last things'' is a scheme of meditation. Its scriptural warrant is one verse of Ecclesiasticus---``In all thy works remember thy last end, and thou shalt never sin''---which counsels remembrance of one's own end and enumerates nothing.\endnote{Ecclesiasticus 7:40, Douay--Rheims from the tracked Challoner text registered in this repository's source library, read 25 July 2026. The verse is the source of the Latin \latin{novissima} behind the English ``last things''; it is a counsel, and the noun in it is one's own end, not a list of four. The character of the fourfold enumeration is treated at length in the companion volume of this series on death, which reaches the same conclusion from its own evidence.} The enumeration built on it---death, judgment, hell, heaven---occurs in none of the acts quoted in this section, and no act that promulgates it has come to this article's notice.\endnote{Bounded, and to be read as bounded. The observation is about the acts actually read for this article: Lateran~IV's \emph{Firmiter}, the Lyons profession, \emph{Ne super his}, \emph{Benedictus Deus}, \emph{Laetentur caeli}, and the two Tridentine texts used in Section~\ref{sec:purgatory}. No search over a magisterial corpus was run and none is claimed. What is asserted is only that these acts define what they define without reference to any fourfold scheme.}

That is not a complaint about the list, which is a good scheme and the frame of this series. It is a warning about what the list cannot show. It has one slot where the acts have two events, and a reader who has met the doctrine through the enumeration arrives at the definitions expecting a single judgment and finds a structure the enumeration has no room to represent. The puzzle of the opening section is therefore sharpened by a devotional ordering that was never built to carry it. Reading the acts in their own sequence dissolves the expectation. It does not dissolve the difficulty.

\subsection{The creed of Lateran IV}

The Fourth Lateran Council opened in 1215 with the constitution \emph{Firmiter}, a creed rather than a canon. Its christological paragraph runs to the end of the world without pausing:

\begin{quote}
\emph{venturus in fine saeculi, iudicaturus vivos et mortuos, et redditurus singulis secundum opera sua, tam reprobis quam electis: qui omnes cum suis propriis resurgent corporibus, quae nunc gestant, ut recipiant secundum opera sua, sive bona fuerint sive mala, illi cum diabolo poenam perpetuam, et isti cum Christo gloriam sempiternam}---about to come at the end of the age, to judge the living and the dead and to render to each according to his works, to reprobate and elect alike: who shall all rise with their own bodies, which they now bear, that they may receive according to their works, whether good or evil, the former with the devil everlasting punishment, the latter with Christ everlasting glory.\endnote{Lateran IV (1215), constitution~1 \emph{Firmiter credimus}, at Denzinger--Sch\"onmetzer 801, read 25 July 2026 in the Latin Enchiridion presented at patristica.net and registered as an exact artifact of that dated web state. The English here is a working gloss of the Latin, not a received translation. The Enchiridion is a compilation and a finding aid to its underlying documents; no printed Denzinger edition and no conciliar manuscript was collated for this article. The Catechism quotes the clause about rising with one's own body at paragraph~999, in its own English.}
\end{quote}

Four things are asserted and one is not. Asserted: a coming at the end of the age; a judgment of living and dead; a rendering to each according to his works; and a general resurrection in each man's \emph{own} body---\emph{quae nunc gestant}, the one he is carrying about at present. Not asserted, anywhere in the sentence: anything whatever about what happens to a man between his death and that day. \emph{Firmiter} is a creed about the last day. It leaves the interval alone.

\subsection{Lyons II, and the sentence that keeps both}

The interval is exactly what the Second Council of Lyons addressed in 1274, in the profession of faith proposed to the Greeks. The relevant paragraphs are three, and they are worth taking as a single movement. Those who die truly penitent in charity before they have satisfied for what they did and left undone are purified after death by purgatorial punishments, and the suffrages of the living---the sacrifices of Masses, prayers, alms, and other offices of piety---avail to relieve them. The souls of those who after baptism contracted no stain of sin at all, and those who contracted it and were purified, whether in the body or out of it, are received into heaven \emph{mox}, straightway. The souls of those who die in mortal sin or in original sin alone descend \emph{mox} into hell, to be punished with unequal punishments. And then, without a paragraph break in sense:

\begin{quote}
\emph{Eadem sacrosancta Ecclesia Romana firmiter credit et firmiter asseverat, quod nihilominus in die iudicii omnes homines ante tribunal Christi cum suis corporibus comparebunt, reddituri de propriis factis rationem}---The same holy Roman Church firmly believes and firmly asserts that \emph{nonetheless} on the day of judgment all men will appear with their bodies before the tribunal of Christ, to render an account of their own deeds.\endnote{Second Council of Lyons (1274), profession of faith of Michael Palaeologus, at DS 856--859; witness and ceiling as in the preceding note. The English is a working gloss. The Latin at DS 859 refers to Romans 14:10f. The article's stress on \emph{nihilominus} is its own reading of the connective and is labelled as project synthesis in the terminal appendix; that the word stands there is a fact of the text.}
\end{quote}

\emph{Nihilominus}. Nonetheless. The council has just said that the dead go at once where they are going, and it immediately adds that this does not cancel the day of judgment. Whoever drafted that paragraph felt the pressure this article is about and refused to release it in either direction.

\subsection{Benedictus Deus}

The most exact of the four texts was occasioned by an embarrassment. In sermons and discussions John~XXII had raised the question whether purified souls separated from their bodies see the divine essence with the vision Paul calls facial before they resume their bodies. On 3~December 1334 he issued \emph{Ne super his}, declaring what his intention had been and confessing that ``the purified souls separated from their bodies are in heaven, the kingdom of heaven, and paradise, and gathered with Christ in the fellowship of the angels, and see God and the divine essence face to face clearly, by the common law, so far as the state and condition of a separated soul allows.''\endnote{John XXII, \emph{Ne super his}, 3~December 1334, at DS 990--991; witness and ceiling as above. English a working gloss of the Latin \emph{Fatemur siquidem et credimus, quod animae purgatae separatae a corporibus sunt in caelo \ldots\ et vident Deum de communi lege ac divinam essentiam facie ad faciem clare, in quantum status et condicio compatitur animae separatae}. The Catechism cites this document beside \emph{Benedictus Deus} at paragraph~1022, footnote~593. This article makes no claim about the character of John~XXII's earlier statements, which it did not read; the declaration is used only to date and locate the controversy that the next document closed.} His successor Benedict~XII closed the question on 29~January 1336. \emph{Benedictus Deus} is worth having in full at the point that matters, because its precision is the whole of its value.

The constitution defines that the souls of all the saints who departed before Christ's passion, and of the apostles, martyrs, confessors, virgins, and other faithful departed after receiving Christ's baptism, \emph{in quibus nihil purgabile fuit}---in whom there was nothing to be purged when they died, nor will be when they die in future---or in whom, if there was something to be purged, \emph{cum post mortem suam fuerint purgatae}, when they have been purged after their death; and the souls of children reborn in the same baptism who die before the use of free will:

\begin{quote}
\emph{mox post mortem suam et purgationem praefatam in illis, qui purgatione huius modi indigebant, etiam ante resumptionem suorum corporum et iudicium generale \ldots\ fuerunt, sunt et erunt in caelo \ldots\ ac \ldots\ viderunt et vident divinam essentiam visione intuitiva et etiam faciali, nulla mediante creatura in ratione obiecti visi se habente, sed divina essentia immediate se nude, clare et aperte eis ostendente}\endnote{Benedict~XII, \emph{Benedictus Deus}, 29~January 1336, at DS 1000; witness and ceiling as above. The English of the central clause is given in the wording of the Catechism, which quotes it at paragraph~1023: souls of the blessed ``have seen and do see the divine essence with an intuitive vision, and even face to face, without the mediation of any creature.'' The remaining renderings in this article's prose are working glosses of the Latin. The registered work record for this constitution already held its identity at catalog level; this article adds no artifact of a printed Denzinger volume and collated none.}
\end{quote}

In English: straightway after their death and the aforesaid purification in those who needed such purification, even \emph{before} the resumption of their bodies and the general judgment, they have been, are, and will be in heaven, and have seen and see the divine essence with an intuitive and even facial vision, no creature intervening as the object seen, but the divine essence showing itself to them immediately, nakedly, clearly, and openly. Souls that die in actual mortal sin descend \emph{mox post mortem suam} to hell.

And then, again, the same connective. Benedict adds: \emph{et quod nihilominus in die iudicii omnes homines `ante tribunal Christi' cum suis corporibus comparebunt, reddituri de factis propriis rationem, `ut referat unusquisque propria corporis, prout gessit, sive bonum sive malum'}---and that nonetheless on the day of judgment all men will appear with their bodies before the tribunal of Christ, to render an account of their own deeds, ``that every one may receive the proper things of the body, according as he hath done, whether it be good or evil.''\endnote{\emph{Benedictus Deus}, at DS 1002; witness and ceiling as above. The internal quotation is 2~Corinthians 5:10, given here in the Douay--Rheims wording checked in the tracked Challoner text registered in this repository (see the note at that verse below). The Latin embeds the Vulgate clause; the constitution's own Latin as delivered by the witness reads \emph{ut referat unusquisque propria corporis, prout gessit, sive bonum sive malum}.} The document that most sharply defines the immediacy of the particular destiny is also the document that most explicitly quotes a general-judgment text of Paul on the far side of it. That is not an oversight. It is the shape of the doctrine.

The Council of Florence repeated the same structure for the Greeks in \emph{Laetentur caeli} (1439), with one addition worth noting: the souls received into heaven ``see clearly God himself, three and one, as he is, yet one more perfectly than another according to the diversity of merits.''\endnote{Council of Florence, bull of union with the Greeks \emph{Laetentur caeli}, 6~July 1439, at DS 1304--1306; witness and ceiling as above. Latin of the addition: \emph{in caelum mox recipi et intueri clare ipsum Deum trinum et unum, sicuti est, pro meritorum tamen diversitate alium alio perfectius}. English a working gloss. The Catechism cites DS 1304--1306 at paragraph~1022, footnote~592, and DS 1304 at paragraph~1031, footnote~604.}

\subsection{What is defined, what is named, and what is open}

Set the four texts side by side and the boundaries come out clean.

\emph{Defined}: that a man's eternal destiny is fixed and effective immediately after death---\emph{mox post mortem}---with purification interposed for those who need it; that the blessed see the divine essence itself, unmediated, before the resurrection; that those who die in mortal sin descend at once to hell; that at the end of the age all men will rise in their own bodies; and that on the day of judgment all will appear before Christ's tribunal to render an account.

\emph{Named, but not defined}: the phrase ``particular judgment.'' It does not occur in any of these acts. What the acts define is the \emph{immediacy of the retribution}, not the existence of a distinctly narrated tribunal at the moment of death. The name is the theologians', and the Catechism has adopted it: ``Each man receives his eternal retribution in his immortal soul at the very moment of his death, in a particular judgment that refers his life to Christ.''\endnote{\emph{Catechism of the Catholic Church} 1022, official English, read 25 July 2026 at the Holy See's archived IntraText presentation and registered as an exact artifact of that page. Its footnotes cite Lyons~II DS 857--858 and Florence DS 1304--1306 and Trent DS 1820 for the purificatory route, \emph{Benedictus Deus} DS 1000--1001 and \emph{Ne super his} DS 990 for the immediate route, and \emph{Benedictus Deus} DS 1002 for immediate damnation. The paragraph closes with the sentence of John of the Cross, ``At the evening of life, we shall be judged on our love.'' The observation that ``particular judgment'' is a theological name for a defined immediacy rather than a separately defined event is this article's, labelled as project synthesis in the terminal appendix; the absence of the phrase from the four acts quoted above is a checkable fact about their wording as delivered by the registered witness.} That is a legitimate and now standard usage, and this article uses it throughout. But a reader should know that the weight sits on \emph{at the very moment of his death}, and that the tribunal imagery around it---an interview, a docket, an advocate, a reading of the record---is theological and devotional elaboration, not defined content.

\emph{Open}: almost everything about the manner of either judgment. Whether the last judgment is conducted vocally or without speech; how long it lasts; where it occurs; the order of its events; what the ``books'' are; whether the account rendered is heard by all or known by all without being heard---none of this is defined, and the medieval authors who discussed it said so. Above all the \emph{when} is not merely undefined but positively withheld, on the Lord's own word: ``But of that day and hour no one knoweth: no, not the angels of heaven, but the Father alone.''\endnote{Matthew 24:36, Douay--Rheims, from the tracked Challoner text of Project Gutenberg eBook 1581 registered in this repository's source library, read 25 July 2026. This is one of two verses quoted in this article that the registered divergence table marks as differing from the settled American edition of 1899, which most printed Douay--Rheims Bibles and the common web presentations reproduce; the American reading is ``But of that day and hour no one knoweth, not the angels of heaven, but the Father alone.'' The difference is punctuation and the emphatic ``no,'' and it does not bear on any claim made here. The other divergent verse is John 5:22, noted at its own place.}

One more text belongs in this list, because it is the most recent and because it states the structural claim of this article as doctrine rather than as inference. In 1979 the Congregation for the Doctrine of the Faith addressed the bishops on what happens between the death of the Christian and the general resurrection, and set out seven points. The fifth: ``In accordance with the Scriptures, the Church looks for `the glorious manifestation of our Lord, Jesus Christ,' believing it to be \emph{distinct and deferred} with respect to the situation of people immediately after death.''\endnote{Sacred Congregation for the Doctrine of the Faith, letter \emph{Recentiores episcoporum Synodi} on certain questions concerning eschatology, 17~May 1979, point~5; Holy See English web text read 25 July 2026 and registered as an exact artifact of that dated state. The internal quotation is from \emph{Dei Verbum} 4. Emphasis added. The letter is an authoritative but non-definitive act of the Roman Curia, approved by John~Paul~II and signed by Cardinal Franjo Seper and Archbishop Jerome Hamer; it is not a conciliar or papal definition, and its weight is stated accordingly wherever it is used here. The Latin published in \emph{Acta Apostolicae Sedis} 71 (1979) governs wording-critical claims and was not consulted; every quotation from this document in this article is identified as the Holy See English.} Distinct and deferred: two judgments, not one judgment described twice.
